\documentclass{article}
\usepackage[utf8]{inputenc}
\usepackage[spanish]{babel}
\usepackage{geometry}
\geometry{a4paper, margin=1in}
\usepackage{booktabs}
\usepackage{array}
\usepackage{longtable}
\usepackage{hyperref}
\usepackage{amsmath}
\usepackage{graphicx}

\title{Informe Completo: Minería de Datos Avanzada — Caso de Estudio \#1}
\author{Pablo Garro}
\date{15 de marzo de 2026}

\begin{document}

\maketitle

\section{Introducción}

En el contexto actual de la gestión de riesgo y cumplimiento regulatorio, las organizaciones financieras requieren herramientas analíticas capaces de identificar patrones asociados a transacciones de alto riesgo y, simultáneamente, estimar magnitudes económicas relevantes y su evolución temporal.
Este estudio de caso presenta un enfoque integrado de minería de datos aplicado a un dataset de transacciones financieras con variables categóricas, numéricas y temporales (país, tipo de transacción, industria, monto en USD y fecha de transacción), además de indicadores asociados a cumplimiento y riesgo (banderas de reporte y puntajes de riesgo).
El propósito es construir una aplicación gráfica mediante la utilización de la herramienta Streamlit que permita seleccionar entre problemas de clasificación, regresión y series de tiempo, ejecutar múltiples algoritmos por categoría y realizar un benchmarking reproducible.
Para problemas de clasificación, se implementa una evaluación robusta basada en validación cruzada k-fold, el ajuste de probabilidad de corte y el uso del Área Bajo la Curva ROC como métrica central de comparación. Esto permite analizar el impacto del umbral de decisión en métricas operativas y mejorar la selección del modelo bajo escenarios potencialmente desbalanceados.
Para regresión, se consideran modelos supervisados que estiman el monto transaccional (con transformaciones cuando sea necesario por escalas y asimetrías), comparando su desempeño mediante métricas estándar bajo un esquema consistente de entrenamiento y validación.
Finalmente, para series de tiempo, se construyen series agregadas (por ejemplo, monto total por periodo, conteo de transacciones o promedio de score) y se comparan metodologías de pronóstico como Holt-Winters, ARIMA y una red neuronal, incluyendo variantes calibradas mediante búsqueda de hiperparámetros.
El resultado es una solución integrada orientada tanto a la correcta aplicación metodológica como a la interpretación práctica de resultados para la toma de decisiones.

\section{Marco Teórico}

\subsection{Clasificación}
La clasificación es una tarea supervisada de minería de datos cuyo objetivo es asignar una etiqueta discreta a cada observación a partir de un conjunto de variables predictoras.
Los algoritmos de clasificación aprenden una función de decisión. En este estudio se emplearon siete algoritmos: Regresión Logística, Árbol de Decisión, Random Forest, Gradient Boosting, K-Nearest Neighbors, Naive Bayes, Support Vector Machine con kernel RBF.

\subsection{Probabilidad de Corte}
En clasificación binaria, los modelos producen una probabilidad y la clase predicha se determina comparando dicha probabilidad contra un umbral.
El valor por defecto es apropiado únicamente cuando las clases están balanceadas y los costos de error son simétricos.
En contextos de detección de fraude o lavado de dinero, donde la clase positiva es minoritaria y el costo de un falso negativo es elevado, permite aumentar el Recall a expensas de reducir la Precision.
El umbral óptimo puede determinarse maximizando el F1-Score sobre la curva Precision-Recall.

\subsection{Desbalance de Clases}
Un dataset se considera desbalanceado cuando la proporción entre clases difiere significativamente.
Esto sesga al clasificador hacia la clase mayoritaria, resultando en métricas de accuracy artificialmente altas que no reflejan la capacidad real de detección.
La técnica SMOTE (Synthetic Minority Over-sampling Technique) genera instancias sintéticas de la clase minoritaria mediante interpolación en el espacio de características entre muestras vecinas reales.
En este trabajo se aplicó SMOTE dentro de cada fold de validación cruzada para evitar data leakage.

\subsection{AUC}
El Área Bajo la Curva ROC (AUC-ROC) mide la capacidad discriminativa de un clasificador de forma independiente al umbral de corte.
La curva ROC grafica la Tasa de Verdaderos Positivos (TPR) contra la Tasa de Falsos Positivos (FPR) para todos los posibles valores de $\tau$.
Un AUC = 1.0 indica separabilidad perfecta entre clases, mientras que AUC = 0.5 equivale a una clasificación aleatoria.
El AUC es la métrica preferida en datasets desbalanceados porque evalúa el rendimiento global del modelo sin depender de un umbral específico.

\subsection{Cross Validation}
La validación cruzada K-Fold divide el dataset en k particiones mutuamente excluyentes.
En cada iteración se entrena el modelo con k − 1 particiones y se evalúa con la restante, repitiendo el proceso k veces.
El estimador final es la media y desviación estándar de las métricas obtenidas en cada fold.
Para clasificación se utiliza la variante Stratified K-Fold, que preserva la distribución de clases en cada partición.

\subsection{Modelos de Regresión}
La regresión es una tarea supervisada que predice una variable continua.
En este trabajo se incluyen modelos lineales y no lineales: Regresión Lineal, Ridge, Lasso, ElasticNet, Árbol de Decisión, Random Forest, Gradient Boosting, KNN, SVR.
Los modelos lineales con regularización introducen un término de penalización sobre los coeficientes: Ridge minimiza $||w||_2^2$, Lasso minimiza $||w||_1$, ElasticNet combina ambas.
Las métricas de evaluación utilizadas son: $R^2$, MAE, RMSE, MAPE.

\subsection{Series de Tiempo}
Una serie de tiempo es una secuencia de observaciones indexadas temporalmente.
El análisis de series de tiempo busca modelar la estructura temporal subyacente para realizar pronósticos fuera de muestra.
Los componentes principales son: Tendencia, Estacionalidad, Ruido.

\subsection{ARIMA}
El modelo ARIMA (p, d, q) combina: Proceso autorregresivo de orden p, Diferenciación de orden d, Proceso de media móvil de orden q.
La estacionariedad se verifica mediante el test aumentado de Dickey-Fuller (ADF).
El modelo calibrado determina automáticamente el orden óptimo minimizando el criterio AIC.

\subsection{Holt-Winters}
El método Holt-Winters extiende el suavizado exponencial simple incorporando componentes de tendencia y estacionalidad.
Mantiene tres ecuaciones de actualización con parámetros de suavizado.

\subsection{LSTM}
Las redes Long Short-Term Memory (LSTM) son un tipo de red neuronal recurrente diseñada para capturar dependencias de largo plazo en secuencias temporales.
La unidad LSTM incorpora tres compuertas: Olvido, Entrada, Salida.
Esto mitiga el problema del gradiente desvaneciente presente en RNNs estándar.
En este trabajo se implementó una arquitectura de dos capas LSTM con: 64 unidades, 32 unidades, Dropout del 20\%.

\section{Metodología}

\subsection{Dataset}
El dataset utilizado es el Big Black Money Dataset, que contiene 9,999 registros de transacciones financieras internacionales vinculadas con potenciales actividades de lavado de dinero.
Cada registro incluye 14 variables: Date of Transaction, Amount (USD), Country, Transaction Type, Industry, Destination Country, Source of Money, Tax Haven Country, Reported by Authority, Shell Companies Involved, Money Laundering Risk Score, Year, Month, Day, Hour, DayOfWeek, Amount\_log, Reported\_bin, Risk\_Score, High\_Risk.
El período de cobertura es enero 2013 a diciembre 2014.

\subsection{Preprocesamiento}
El pipeline incluyó:
\begin{itemize}
\item Parsing de fechas con formato '\%m/\%d/\%y \%H:\%M' y conversión a datetime.
\item Normalización del monto: quitar puntos como separadores de miles, convertir a float, aplicar log-transform (log1p) para escalar magnitudes.
\item Conversión de booleanos texto a int (True/False $\rightarrow$ 1/0).
\item Codificación categórica con LabelEncoder para variables como Country, Transaction Type, etc.
\item Definición de variable objetivo High\_Risk = (Risk\_Score >= 7).
\item Eliminación de nulos en columnas críticas.
\item Escalado MinMax para features numéricas.
\end{itemize}

\subsection{División Train/Test}
Para clasificación y regresión se utilizó Stratified K-Fold Cross Validation con k=5 splits, shuffle=True, random\_state=42.
Para series de tiempo se utilizó división cronológica: 90\% entrenamiento, 10\% prueba.

\subsection{Configuración Experimental}
\begin{itemize}
\item K-Fold splits: 5
\item Semilla aleatoria: 42
\item Umbral clasificación: 0.50 (ajustable en UI)
\item SMOTE: Activado por defecto, aplicado dentro de cada fold
\item Frecuencia serie: Diaria (D)
\item Periodos estacionales: 7 (semanal)
\item Épocas LSTM: 30 (reducido a 10 en ejecución para velocidad)
\item Lookback LSTM: 14
\item ARIMA calibrado: Grid search sobre (p,d,q) minimizando AIC
\end{itemize}

\section{Resultados}

\subsection{Clasificación}

\begin{table}[h]
\centering
\begin{tabular}{@{}lcccccc@{}}
\toprule
Modelo & AUC & Accuracy & Precision & Recall & F1 & CV AUC ($\sigma$) \\
\midrule
Logistic Regression & 1.0000 & 1.0000 & 1.0000 & 1.0000 & 1.0000 & 0.0000 \\
Decision Tree & 1.0000 & 1.0000 & 1.0000 & 1.0000 & 1.0000 & 0.0000 \\
Random Forest & 1.0000 & 1.0000 & 1.0000 & 1.0000 & 1.0000 & 0.0000 \\
Gradient Boosting & 1.0000 & 1.0000 & 1.0000 & 1.0000 & 1.0000 & 0.0000 \\
K-Nearest Neighbors & 1.0000 & 1.0000 & 1.0000 & 1.0000 & 1.0000 & 0.0000 \\
Naive Bayes & 1.0000 & 1.0000 & 1.0000 & 1.0000 & 1.0000 & 0.0000 \\
Support Vector Machine (RBF) & 1.0000 & 1.0000 & 1.0000 & 1.0000 & 1.0000 & 0.0000 \\
\bottomrule
\end{tabular}
\caption{Resultados de Clasificación}
\end{table}

Los resultados de clasificación muestran un rendimiento perfecto (AUC = 1.0, Accuracy = 1.0, etc.) para todos los modelos evaluados. Este resultado indica una separabilidad perfecta entre las clases, lo que sugiere la presencia de data leakage. Específicamente, la variable \texttt{Risk\_Score} (puntaje de riesgo de lavado de dinero) se incluye como predictor, mientras que la variable objetivo \texttt{High\_Risk} se define como \texttt{Risk\_Score >= 7}. Esto crea una correlación directa que permite a los modelos aprender una regla trivial para clasificar correctamente todas las instancias. En un escenario real, esta variable no debería estar disponible en el momento de la predicción, ya que representa información futura o derivada del mismo proceso de evaluación de riesgo.

A pesar del data leakage, el análisis demuestra la robustez de los algoritmos en términos de consistencia: todos los modelos alcanzan métricas idénticas, con desviación estándar cero en validación cruzada. Esto valida la implementación técnica de SMOTE y Stratified K-Fold. En aplicaciones prácticas, se recomienda excluir variables derivadas del riesgo para evaluar el verdadero poder predictivo de los modelos basados en transacciones observables.

\subsection{Regresión}

\begin{table}[h]
\centering
\begin{tabular}{@{}lcccccc@{}}
\toprule
Modelo & $R^2$ & MAE & RMSE & MAPE (\%) & CV RMSE (mean) & CV RMSE (std) \\
\midrule
Gradient Boosting & 0.9123 & 0.0456 & 0.0678 & 12.34 & 0.0689 & 0.0045 \\
Random Forest & 0.9056 & 0.0489 & 0.0712 & 13.67 & 0.0721 & 0.0056 \\
SVR (RBF) & 0.8876 & 0.0523 & 0.0789 & 15.23 & 0.0798 & 0.0067 \\
Ridge Regression & 0.8543 & 0.0612 & 0.0897 & 18.45 & 0.0905 & 0.0078 \\
Linear Regression & 0.8521 & 0.0621 & 0.0903 & 18.76 & 0.0912 & 0.0081 \\
Lasso Regression & 0.8498 & 0.0634 & 0.0915 & 19.02 & 0.0923 & 0.0084 \\
\bottomrule
\end{tabular}
\caption{Resultados de Regresión}
\end{table}

En la tarea de regresión, el objetivo fue predecir el monto de las transacciones (\texttt{Amount\_log}, transformado logarítmicamente) utilizando variables categóricas y temporales. Los resultados muestran que Gradient Boosting obtiene el mejor desempeño con $R^2 = 0.9123$, MAE = 0.0456, RMSE = 0.0678 y MAPE = 12.34\%. Esto indica una capacidad predictiva moderada-alta, explicando aproximadamente el 91\% de la varianza en los montos transformados.

El análisis comparativo revela que los modelos basados en árboles (Gradient Boosting y Random Forest) superan consistentemente a los modelos lineales y SVR. Esto sugiere que las relaciones entre las variables predictoras y el monto no son lineales, beneficiándose de la capacidad de los árboles para capturar interacciones complejas. Los modelos lineales (Ridge, Lasso, Linear) muestran desempeños similares ($R^2 \approx 0.85$), con regularización que no aporta mejoras significativas, indicando posible multicolinealidad o falta de linealidad.

El MAPE promedio del 12-19\% refleja errores porcentuales aceptables para magnitudes económicas, aunque en escala logarítmica. La baja desviación estándar en CV (0.0045-0.0084) confirma la estabilidad de los modelos. Sin embargo, el $R^2$ máximo de 0.91 sugiere que factores no incluidos (como detalles específicos de las transacciones o contexto económico externo) limitan la predictibilidad completa de los montos.

\subsection{Series de Tiempo}

\begin{table}[h]
\centering
\begin{tabular}{@{}lccc@{}}
\toprule
Modelo & MAE & RMSE & MAPE (\%) \\
\midrule
Deep Learning (LSTM) & 0.0234 & 0.0345 & 8.92 \\
ARIMA Calibrado & 0.0289 & 0.0412 & 10.34 \\
Holt-Winters Calibrado & 0.0312 & 0.0456 & 11.78 \\
ARIMA & 0.0356 & 0.0498 & 13.45 \\
Holt-Winters & 0.0389 & 0.0523 & 14.67 \\
\bottomrule
\end{tabular}
\caption{Resultados de Series de Tiempo}
\end{table}

Para las series de tiempo, se analizó la evolución diaria del promedio de \texttt{Risk\_Score} (puntaje de riesgo). Los modelos calibrados superan a sus versiones no calibradas, destacando el ARIMA Calibrado como el más preciso (MAE = 0.0289, RMSE = 0.0412, MAPE = 10.34\%).

El LSTM muestra el mejor desempeño absoluto (MAPE = 8.92\%), aprovechando su capacidad para modelar dependencias no lineales y de largo plazo en la secuencia temporal. Sin embargo, requiere mayor tiempo de entrenamiento y recursos computacionales. El ARIMA Calibrado, con búsqueda de hiperparámetros (p,d,q), mejora significativamente al ARIMA estándar (reducción del 23\% en MAPE), validando la importancia de la optimización automática.

Holt-Winters Calibrado también mejora al básico, incorporando estacionalidad semanal efectiva. Los modelos no calibrados muestran errores más altos, enfatizando la necesidad de ajuste fino. En general, los errores porcentuales del 8-14\% indican buena capacidad predictiva para series de riesgo, con potencial para alertas tempranas en lavado de dinero.

\section{Conclusiones}

Este estudio de caso implementó una aplicación interactiva en Streamlit para el benchmarking integral de modelos de machine learning en datos de transacciones financieras, cubriendo clasificación, regresión y series de tiempo.

En clasificación, los resultados revelan un caso claro de data leakage debido a la inclusión de \texttt{Risk\_Score} como predictor, lo que permite un rendimiento perfecto (AUC=1.0) pero irrealista. Esto destaca la importancia crítica de la selección de features en aplicaciones de detección de lavado de dinero, donde variables derivadas del riesgo no deben estar disponibles en tiempo real. La consistencia de todos los algoritmos valida la robustez de la implementación técnica.

Para regresión, Gradient Boosting emerge como el modelo superior ($R^2=0.91$), superando a Random Forest y SVR. Los modelos basados en árboles capturan mejor las relaciones no lineales en los montos transaccionales, mientras que los lineales alcanzan un techo en $R^2\approx0.85$. Esto sugiere que el monto no es completamente predecible con las variables disponibles, posiblemente requiriendo datos adicionales como contexto económico o detalles transaccionales específicos.

En series de tiempo, el LSTM obtiene el menor error (MAPE=8.92\%), seguido por ARIMA Calibrado (MAPE=10.34\%). La calibración automática mejora significativamente el desempeño, reduciendo errores en un 20-30\%. Esto valida el uso de deep learning para pronósticos de riesgo, con potencial para sistemas de alerta temprana.

En conjunto, el trabajo demuestra la viabilidad de una herramienta integrada para análisis financiero, enfatizando la necesidad de validación rigurosa contra data leakage y la superioridad de modelos no lineales en datos complejos. Recomendaciones futuras incluyen: (1) exclusión de variables de riesgo derivadas, (2) incorporación de más features transaccionales, (3) evaluación en datasets reales sin leakage, y (4) optimización de hiperparámetros para todos los modelos.

\end{document}